\subsection{Research Gap and Proposed Approach}

The unresolved research gap is the absence of a jointly evaluated path from a telemetry alarm to an engineering inspection lead. UAV anomaly-management and fault-diagnosis studies primarily evaluate whether abnormal behavior can be detected or assigned to a fault category \cite{fourlas2021survey,tan2026anomaly,wei2026aerotsboost}. Explainable machine-learning studies instead evaluate how model predictions can be attributed to input features \cite{lundberg2017shap,lundberg2020trees}, while software fault-localization studies rank code entities using program-level evidence \cite{wong2016survey}. These lines address different parts of the diagnostic process, but their outputs and ground truths are not directly interchangeable. Evaluating them separately therefore leaves unanswered whether one abnormal flight can be followed consistently from detection, through fault-domain diagnosis and telemetry evidence, to a defensible set of PX4 module inspection leads.

A second gap concerns traceability across representation levels. Temporal-statistical models operate on descriptors such as windowed means, extrema, variations, or slopes, whereas a flight-control engineer reasons about physical telemetry channels, logged topics, and functional subsystems. Feature attributions can quantify the contribution of individual model inputs, but they do not inherently recover this engineering hierarchy. Moreover, a PX4 module relationship cannot be inferred from a feature name alone: it depends on the firmware version and on the uORB topics published or subscribed by each module \cite{meier2015px4,px4uorbdocs2026}. Without an explicit and quantitatively conserved mapping from statistical feature to telemetry channel, topic, subsystem, and module, feature importance remains difficult to translate into reproducible inspection priorities.

A third gap is the validation of software-module rankings. Static topic dependencies show architectural connectivity, not execution, defect, or causality. Software fault-localization research evaluates suspiciousness rankings against known faulty entities, and prior work has shown that conclusions can depend strongly on the faults and evaluation protocol used \cite{wong2016survey,pearson2017evaluating}. Controlled source mutations can provide explicit module ground truth, although mutation evidence must itself be interpreted with care because artificial and real faults are not equivalent \cite{papadakis2019mutation}. For a telemetry-driven cascade, validation must additionally preserve the detection gate: an onset-conditioned ranking may reveal useful downstream evidence when an anomalous interval is supplied, but it does not constitute an end-to-end localization result if the runtime detector never activates.

To address these gaps, this paper develops HS-AeroTS as a hierarchical diagnostic evidence framework. The earlier ``FL'' suffix is removed because the evaluated system does not establish fault localization. A first-stage detector scores telemetry windows as normal or anomalous, and a second-stage classifier assigns detected anomalies to External Position, Global Position, Altitude, or Mechanical/Electrical domains. TreeSHAP evidence is then aggregated through an explicit feature--channel--topic--subsystem hierarchy, after which a static PX4 uORB graph reconstructed at the majority firmware commit propagates topic evidence from matching-commit windows to software-module suspects. The framework keeps each evidence level separately testable: detection and fault-domain diagnosis are evaluated on real-flight labels, explanation quality is assessed through conservation, masking, and semantic-mapping sensitivity, and module ranking is examined with controlled mutation replay under both onset-conditioned and detector-gated protocols. This separation produces a connected diagnostic chain while preventing architecture-aware associations from being presented as source-code root causes.
